%% generated by gen_tables.py - RE-ANALYSIS 2026
\begin{table}[htbp]
\centering
\small
\caption{Structural load cases and support scenarios (idealised boundary conditions)}
\label{tab:supports}
\begin{tabular}{>{\raggedright\arraybackslash}p{0.20\linewidth}>{\raggedright\arraybackslash}p{0.52\linewidth}>{\raggedright\arraybackslash}p{0.18\linewidth}}
\toprule
Case & Constraint definition & Role \\
\midrule
LC1 free expansion & 3 outer-ring nodes on the inlet face ($\theta$ = 0/120/240\textdegree), $U_\theta = U_z = 0$ in the duct cylindrical system; statically determinate & gradient stress only \\
LC2 = S1 & $U_z = 0$ on both complete end faces; $U_\theta = 0$ at 3 mid-span outer nodes; radial free; ends free to sway, end rotation held (guided column, $K = 1$) & official baseline \\
S2 & $U_z = U_\theta = 0$ on every node of both end faces (clamped--clamped, $K = 0.5$) & sensitivity \\
S3 & inlet face $U_z = U_\theta = 0$ (clamped); outlet on a deformable remote point on the axis, $U = 0$, rotations free (pinned, $K = 0.699$) & sensitivity \\
\bottomrule
\end{tabular}
\par\smallskip\parbox{\linewidth}{\footnotesize\raggedright Source: \texttt{MASTER\_PROJECT\_DATA.csv} rows M019--M022. None of these represents a defined real flange; the real end restraint is undefined (open task T-034).}
\end{table}
